\documentclass[10pt]{article}
\usepackage[T1]{fontenc}
\usepackage[utf8]{inputenc}
\usepackage[french]{babel}
\usepackage{amsmath,amssymb,mathrsfs,mathpazo}
\usepackage{geometry}
\geometry{a4paper,margin=2cm}
\usepackage{graphicx,array,tabularx,tabularray,colortbl,multirow,multicol,booktabs,enumitem,pifont,ulem,fancybox,diagbox,cancel}
\usepackage{pstricks,pstricks-add,pst-plot,pst-tree,pst-node,pst-text,pst-func,pst-math,pst-bspline,pst-3dplot,pst-eucl,pst-solides3d}
\let\pscalculate\undefined
\usepackage{pst-bezier}
\usepackage{tikz,pgfplots,tkz-tab}
\usetikzlibrary{arrows,arrows.meta,patterns,calc,positioning,trees,decorations.markings,shapes.geometric,shapes.misc,shapes.arrows,intersections,decorations.pathmorphing,decorations.pathreplacing,backgrounds,shadows,babel}
\pgfplotsset{compat=1.18}
\usepgfplotslibrary{fillbetween}
\usepackage[np]{numprint}
\usepackage{listings,makecell,pas-tableur,twoopt}
\let\barre\undefined
\usepackage[pstricks]{tablvar}
\definecolor{darkgreen}{RGB}{0,100,0}
\providecommand{\red}{\color{red}}
\providecommand{\blue}{\color{blue}}
\providecommand{\green}{\color{green}}
\providecommand{\black}{\color{black}}
\providecommand{\white}{\color{white}}
\providecommand{\gray}{\color{gray}}
\providecommand{\cyan}{\color{cyan}}
\providecommand{\magenta}{\color{magenta}}
\providecommand{\yellow}{\color{yellow}}
\providecommand{\violet}{\color{violet}}
\providecommand{\orange}{\color{orange}}
\providecommand{\euro}{\texteuro}
\providecommand{\EUR}{\texteuro}
\providecommand{\ieme}{\textsuperscript{e}}
\providecommand{\degres}{\ensuremath{^\circ}}
\providecommand{\degree}{\ensuremath{^\circ}}
\providecommand{\vv}[1]{\overrightarrow{#1}}
\providecommand{\Coord}[1]{#1}
\providecommand{\decosix}{}
\providecommand{\starredbullet}{\ensuremath{\bullet}}
\pagestyle{empty}
\setlength{\parindent}{0pt}
\usepackage{adjustbox}\DeclareUnicodeCharacter{2236}{:}\DeclareUnicodeCharacter{2190}{\ensuremath{\leftarrow}}
\providecommand{\R}{\mathbb{R}}
\providecommand{\N}{\mathbb{N}}
\providecommand{\Z}{\mathbb{Z}}
\providecommand{\Q}{\mathbb{Q}}
\providecommand{\D}{\mathbb{D}}
\providecommand{\C}{\mathbb{C}}
\providecommand{\e}{{\mathrm{e}}}
\providecommand{\vect}[1]{\overrightarrow{#1}}
\providecommand{\vectt}[1]{\overrightarrow{\mathrm{#1}}}
\providecommand{\vvt}[1]{\overrightarrow{\mathrm{#1}}}
\providecommand{\np}[2][]{{#2~#1}}
\providecommand{\vv}[1]{{\overrightarrow{#1}}}
\renewcommand{\d}{{\mathrm{d}}}
\providecommand{\degres}{{^\circ}}
\providecommand{\degre}{{^\circ}}
\providecommand{\blue}{\color{blue}}
\providecommand{\red}{\color{red}}
\providecommand{\enskip}{\quad}
\providecommand{\llbracket}{\lbrack\!\lbrack}
\providecommand{\rrbracket}{\rbrack\!\rbrack}
\providecommand{\euro}{{\text{€}}}
\providecommand{\no}{{\text{n°}}}
\providecommand{\interval}[2]{{[#1;#2]}}
\providecommand{\intervFO}[2]{{[#1;#2[}}
\providecommand{\dotfill}{\ldots}

\providecommand{\ds}{\displaystyle}
\providecommand{\textding}[1]{#1}
\providecommand{\Oij}{\(\mathrm{O};\vec{\imath},\vec{\jmath}\)}
\providecommand{\Oijk}{\(\mathrm{O};\vec{\imath},\vec{\jmath},\vec{k}\)}
\providecommand{\Ouv}{\(\mathrm{O};\vec{u},\vec{v}\)}
\providecommand{\intervFF}[2]{\ensuremath{[#1;#2]}}
\renewcommand{\euro}{\ifmmode\text{\texteuro}\else\texteuro\fi}
\DeclareUnicodeCharacter{00B0}{\ifmmode^{\circ}\else\textdegree\fi}
\DeclareUnicodeCharacter{20AC}{\ifmmode\text{\texteuro}\else\texteuro\fi}
% Énoncés : Éducation nationale. Transcription et corrigés : APMEP, Denis Vergès et contributeurs. Découpage et classement : CoopMaths, Sébastien Lozano, Éric Elter et contributeurs. Adaptation de la présentation : Maths Explorer.
% CC BY-SA : https://coopmaths.fr/www/about/
% https://coopmaths.fr/alea/static/bac/2023/tex/bac_2023_03_sujet1_etrangers_groupe2_4_cor.tex
\begin{document}
\noindent\textit{APMEP / CoopMaths. CC BY-SA. Présentation adaptée par Maths Explorer.}\par\medskip
\medskip

Dans l'espace muni d'un repère orthonormé \Oijk, on considère les
points
\begin{center}A$(-1~;~-3~;~2)$,\qquad B$(3~;~-2~;~6)$ \quad et \quad  C$(1~;~2~;~-4)$.\end{center}

\smallskip

\begin{enumerate}
\item 
On a $\vect{\text{AB}}\begin{pmatrix}4\\1\\4\end{pmatrix}$ et $\vect{\text{AC}}\begin{pmatrix}2\\5\\-6\end{pmatrix}$ : ces vecteurs ne sont pas colinéaires, donc les points A, B et C ne sont pas alignés : ils définissent donc un plan $\mathcal{P}$.
\item 
	\begin{enumerate}
		\item 
On a :

$\bullet~~$$\vect{n} \cdot \vect{\text{AB}} = 4 \times 13 + 1 \times (- 16) + 4 \times (-9) = 52 - 16 - 36 = 52 - 52 = 0$ ;

$\bullet~~$$\vect{n} \cdot \vect{\text{AC}} = 2 \times 13 + 5 \times (- 13) + (- 6 \times (-9) = 26 - 65 + 54  = 52 - 52 = 0$.

Le vecteur $\vect{n}$ est donc orthogonal à deux vecteurs non colinéaires du plan $\mathcal{P}$ : il est normal à ce plan.
		\item 
On sait que $M(x~;~y~;~z) \in \mathcal{P} \iff ax + by + cz + d = 0$ et que $a, b $ et $c$ sont les composantes d'un vecteur normal à ce plan donc par exemple $\vect{n}$.
		
On a donc $M(x~;~y~;~z) \in \mathcal{P} \iff 13x - 16y -9z + d = 0$.
		
par exemple C$(1~;~2~;~- 4) \in \mathcal{P} \iff 13\times 1 + 2\times  (-16) -9 \times (- 4) + d = 0   + d = 0 \iff 13 - 32 + 36 + d = 0 \iff 17 + d = 0 \iff d = - 17$.
		
On a donc $M(x~;~y~;~z) \in \mathcal{P} \iff 13x - 16y - 9z - 17 = 0$.
	\end{enumerate}
\end{enumerate}

On note $\mathcal{D}$ la droite passant par le point F$(15~;~-16~;~-8)$ et orthogonale au plan $\mathcal{P}$.
\begin{enumerate}[resume]
\item 
La droite $\mathcal{D}$ étant orthogonale au plan $\mathcal{P}$ a pour vecteur directeur le vecteur $\vect{n}$ normal au plan $\mathcal{P}$.

On a donc $M(x~;~y~;~z) \in \mathcal{D} \iff \vect{\text{F}M} = t \vect{n}, \quad t \in \R$, soit $\left\{\begin{array}{l c l}
x - 15&=&13t\\
y + 16&=&- 16t\\
z + 8&=&- 9t
\end{array}\right., \quad t \in \R$
\item 
Si E$(x~;~y~;~z)$ est commun à la droite $\mathcal{D}$ et au plan $\mathcal{P}$, ses coordonnées vérifient les équations de la droite et celle du plan donc le système :

$\left\{\begin{array}{l c l}
x - 15&=&13t\\
y + 16&=&- 16t\\
z + 8&=&- 9t\\
13x - 16y - 9z - 17 &=& 0
\end{array}\right., \quad t \in \R$.

En remplaçant $x, y$ et $z$ par leurs valeurs en fonction de $t$ dans l'équation du plan, on obtient :

$13(13t + 15) - 16(16t - 16)- 9(- 9t - 8) - 17 = 0 \iff 169t + 195 + 256t + 256 +81t + 72 - 17 = 0 \iff 506t + 506 = 0 \iff t = - 1$.

En reportant cette valeur dans les équations paramétriques de la droite $\mathcal{D}$, on obtient  :

$\left\{\begin{array}{l c l}
x - 15&=&-13 \\
y + 16&=&16\\
z + 8&=&9\\
\end{array}\right.\iff \left\{\begin{array}{l c l}
x &=&2 \\
y &=&0\\
z &=&1\\
\end{array}\right.$

Donc E(2~;~0~;~1).

\item 
F et E appartiennent à la droite $\mathcal{D}$ perpendiculaire au plan $\mathcal{P}$, donc la distance du point F au plan est égale à FE.

Or FE$^2 = (- 13)^2 + 16^2 + 9^2 = 169 + 256 + 81 = 506$. D'où FE $ = \sqrt{506}$.
\item 
Comme précédemment si $M \in \mathcal{D}$ perpendiculaire au plan $\mathcal{P}$ la distance de ce point au plan $\mathcal{P}$ est $M$E.

\begin{center}
\psset{unit=1cm}
\begin{pspicture}(0,-1)(8,1)
\psline(8,0)
\psdots(1,0)(3,0)(5,0)(7,0)
\uput[d](1,0){F} \uput[d](3,0){$M_1$} \uput[d](5,0){E} \uput[d](7,0){$M_2$}
\multido{\n=1.9+2.0,\na=2.1+2.0}{3}{\psline(\n,-0.2)(\na,0.2)}
\end{pspicture}
\end{center}

$\bullet~~$Premier point répondant à la question : $M_1$ tel que $\vect{\text{E}M_1} = \dfrac12 \vect{\text{EF}}$.

De $\vect{\text{EF}}\begin{pmatrix}13\\-16\\-9\end{pmatrix}$, on déduit que $\dfrac12 \vect{\text{EF}} = \vect{\text{E}M_1}\begin{pmatrix}6,5\\-8\\-4,5\end{pmatrix}$, soit 

$\left\{\begin{array}{l c l}
x_1 - 2&=&6,5\\
y_1 - 0&=&- 8\\
z - 1&=&-4,5
\end{array}\right. \iff \left\{\begin{array}{l c l}
x_1 &=&8,5\\
y_1 &=&- 8\\
z &=&-3,5
\end{array}\right.$

$\bullet~~$Deuxième point répondant à la question : $M_2$ tel que $\vect{\text{E}M_2} = -\dfrac12 \vect{\text{EF}}$.

De $\vect{\text{EF}}\begin{pmatrix}13\\-16\\-9\end{pmatrix}$, on déduit que $-\dfrac12 \vect{\text{EF}} = \vect{\text{E}M_2}\begin{pmatrix}-6,5\\8\\4,5\end{pmatrix}$, soit 

$\left\{\begin{array}{l c l}
x_2 - 2&=&-6,5\\
y_2 - 0&=& 8\\
z_2 - 1&=&4,5
\end{array}\right. \iff \left\{\begin{array}{l c l}
x_2 &=&- 4,5\\
y_2 &=& 8\\
z_2 &=&5,5
\end{array}\right.$

Donc $M_1(8,5~;~-8~;~- 3,5)$ et $M_2(-4,5~;~ 8~;~5,5)$.
\end{enumerate}

\end{document}
